\documentclass[../../../include/open-logic-section]{subfiles}

\begin{document}

\olfileid{sth}{cardinals}{zfc}
\olsection{$\ZFC$: एक महत्त्वाचा टप्पा}

सुक्रमण जोडल्यावर आपण उपपत्तीचा अंतिम महत्त्वाचा टप्पा गाठला आहे. आता~$\ZFC$ साठी आवश्यक असलेली सर्व स्वयंसिद्धके आपल्याकडे आहेत. सविस्तर यादी अशी:

\begin{defn}
$\ZFC$ उपपत्तीत पुढील स्वयंसिद्धके आहेत: विस्तारात्मकता, संयोग, जोड्या, घातसंच, अनंतता, पायाभूतता, सुक्रमण आणि विभक्तीकरण व प्रतिस्थापन या रूपबंधांतील सर्व स्वयंसिद्धके. दुसऱ्या शब्दांत, $\ZF$ मध्ये सुक्रमण जोडले की $\ZFC$ मिळते.
\end{defn}

$\ZFC$ म्हणजे \emph{झर्मेलो--फ्रेंकेल} संच उपपत्ती, \emph{निवडीच्या स्वयंसिद्धकासह}. आपण जोडलेल्या स्वयंसिद्धकाला “निवड” नव्हे तर “सुक्रमण” म्हटले, म्हणून हे थोडे विचित्र वाटू शकते. पण पुढे आपण {निवडीचे स्वयंसिद्धक} मांडल्यावर, सुक्रमण आणि निवड ही $\ZF$ गृहीत धरल्यास समतुल्य आहेत हे दिसेल (\olref[choice][woproblem]{thmwochoice} पाहा). त्यामुळे यांपैकी कोणते “मूलभूत” स्वयंसिद्धक मानायचे, याने फरक पडत नाही. आणि “$\ZFC$” हे नाव साहित्यात सर्वमान्य आहे.

\end{document}
